\section{O Problema da Calendarização Universitária}
\label{sec:2.1}

O \textit{University Course Timetabling Problem} (UCTP) consiste em alocar eventos acadêmicos, como aulas, laboratórios ou avaliações, a intervalos de tempo e, em sua formulação mais ampla, a espaços físicos disponíveis, respeitando simultaneamente disponibilidades docentes, restrições curriculares, limitações institucionais e critérios pedagógicos. \citeonline[p.~1]{ceschia2023} definem o problema como a tarefa de atribuir reuniões entre professores e alunos a dias, intervalos de tempo e salas de aula, ressaltando que, na prática, cada instituição possui suas próprias regras, convenções e especificidades, o que torna cada instância do problema fortemente dependente do contexto institucional.

Do ponto de vista da teoria da complexidade computacional, o UCTP é reconhecido por sua alta complexidade combinatória, o que dificulta sua resolução por métodos exatos em instâncias gerais \cite[p.~1]{premananda2024}. À medida que diferentes variáveis e restrições são consideradas simultaneamente, o espaço de busca cresce de forma significativa, tornando necessária a adoção de estratégias aproximativas capazes de encontrar soluções satisfatórias em tempo computacional aceitável.

As restrições do problema são geralmente organizadas em duas categorias. As restrições rígidas, conhecidas como \textit{hard constraints}, são mandatórias, de modo que sua violação compromete a validade da grade produzida. Um exemplo típico é a proibição de alocar um mesmo professor em dois horários simultâneos. As restrições flexíveis, conhecidas como \textit{soft constraints}, não invalidam necessariamente a solução, mas influenciam sua qualidade, como a redução de janelas ociosas para alunos, a distribuição equilibrada da carga horária ou o atendimento a preferências de turno dos professores \cite{weare1995}. Em processos manuais, a montagem da grade tende a priorizar o atendimento das restrições rígidas, enquanto as restrições flexíveis podem ser parcialmente negligenciadas devido à complexidade do planejamento. Essa limitação justifica o uso de abordagens heurísticas e meta-heurísticas para apoiar a geração de soluções viáveis em tempo computacional aceitável.
